\section{Data}
\label{sec:data}

\subsection{Source and vintage}

All data are from the Human Mortality Database \citep{hmd2026}, bulk
``zipped data files'' download, files \texttt{Deaths\_1x1},
\texttt{Exposures\_1x1}, \texttt{Mx\_1x1} and the period life tables, last
modified 15~June~2026, release DOI 10.4054/HMD.Countries.20260615, and
constructed under Methods Protocol version~6, revision of 5~August~2025
\citep{wilmoth2025methods}. HMD has issued a citable DOI per release since
September 2025; the DOI, the download date (2026-08-25) and the SHA-256
checksums of every input file in \texttt{data/MANIFEST.sha256} together pin
the exact bytes. The data themselves are not redistributed, in keeping with
the HMD user agreement, and any table in this paper can be reproduced from
a fresh download that matches the manifest. The bulk deaths and exposures
files each carry fifty distinct population codes (HMD's \texttt{PopName}),
the count recorded in the manifest header: forty-one countries or areas,
plus nine further codes for the overlapping sub-population and civilian
variants that HMD publishes for Germany (two), France (one), the United
Kingdom (four) and New Zealand (two). This is the first HMD vintage in which
twenty of them carry \emph{final} single-age annual data through 2024,
which is what makes the shift regime (Section~\ref{sec:design}) possible at
$h=5$; one (DNK) already reaches 2025, which the shift regime does not use.

Only the deaths and exposures enter the models. The \texttt{Mx\_1x1} file is
used once, to validate the parser: our $D/E$ reproduces the published $m_x$
exactly at spot checks. The period life tables are used once, for validation
gate~3 (Section~\ref{sec:design-gates}).

\subsection{Cells, ages and exposure}

A cell is a (population, sex, calendar year, single age) Lexis square with
death count $D_{x,t}$ and central exposure $E_{x,t}$. HMD reports ages
$0,\dots,109$ and an open group $110{+}$. We model ages $0$--$99$ and drop
the rest, for three reasons. First, the published life tables smooth death
rates at the oldest ages (a Kannisto fit above age~80 under Protocol~v6), so
the old-age surface is partly model output rather than observation. Second,
the open group needs a closure assumption that no forecasting family in the
grid shares. Third, exposure vanishes: read with the study's own parser
(\texttt{mortcal.data.hmd.read\_bulk\_1x1}), the vintage contains 24{,}083
population--year--age rows with $E_{x,t}=0$ for at least one sex (15{,}494
for both sexes), almost all at age~90 or above and concentrated in the
small populations; a Poisson model with offset $\log E_{x,t}$ is undefined
there. Sexes are modelled separately throughout; the \emph{Total} column is
never used.

Inside $0$--$99$ a zero-exposure cell is a \emph{structural} zero, not a
missing value: in every such cell in the shift training panel HMD's
1~January population count is zero at both corners of the Lexis square, so
nobody was alive at that age and $D_{x,t}=0$ as well. Under the Poisson
likelihood with log-exposure offset such a cell has mean $E\exp(\eta)=0$
and observation $0$: its log-likelihood contribution, score and Fisher
information are all identically zero, so the maximum-likelihood estimate
with the cell included equals the estimate under indicator weights
\emph{exactly}; only the numerical evaluation of $0\log 0$ fails. Every
family therefore carries the cell-weight matrix $W_{x,t}=\ind\{E_{x,t}>0\}$
(the \texttt{wxt} convention of \pkg{StMoMo}), which Addendum~3~\S1 of the
pre-registration records as the numerically correct evaluation of the
registered objective rather than as a deviation; the weighted code paths
reproduce the unweighted fits bit-identically on panels without zeros,
which is tested. The shift training panel contains 99 such cells, all at
ages 95--99, in seven population--sex series (CHE, FIN, ISL, LUX and SWE
males, CHE and ISL females), every one before 1970. In a \emph{test}
window a zero-exposure cell has no observation to score; the scorer masks
an age only if every horizon's observation is finite, records the number of
ages scored, and computes life-table functionals on the maximal contiguous
scored age range from age~0 on both the predictive and the observed side
(Addendum~3~\S3). This binds only in the placebo regime, at ages 98--99, in
a handful of Swiss, Finnish and Icelandic cells. Zero-\emph{death} cells
with positive exposure are ordinary observations under the Poisson
likelihood and remain in every fit.

\subsection{Fractional deaths, half-count rates and the rounding convention}
\label{sec:data-rounding}

HMD death counts are not integers. Deaths are allocated to Lexis triangles
and then re-aggregated, and deaths of unknown age are redistributed, so a
cell may record $3842.31$ deaths; this is the norm, not an anomaly. It has
two consequences that the pre-registration and its addenda fix in advance.

First, a Poisson log score evaluated at a non-integer count is undefined.
We score the Poisson predictive at the death count rounded to the nearest
integer with halves rounded \emph{up}, $\tilde D_{x,t} = \lfloor D_{x,t} +
\tfrac12 \rfloor$, and report as a registered sensitivity the CRPS of the
predictive count distribution against the unrounded $D_{x,t}$, which needs
no rounding. The half-up rule is implemented as written; IEEE
round-half-to-even was rejected because a small share of HMD cells land
exactly on a half and its direction would then depend on the parity of the
integer part, a cell-dependent perturbation no reader would infer from
``nearest integer''. Nine of the twenty shift populations carry
non-integer counts, two of them (JPN, USA) in essentially every cell; the
other eleven are exactly integer, and the per-population non-integer share
is reported so that the rounding's reach is visible.

Second, a cell with zero observed deaths has no finite log rate. Rate-scale
quantities---observed and predictive alike---use the half-count continuity
correction of demographic practice,
\begin{equation}
  y_{x,t} = \log\frac{\max(\tilde D_{x,t},\,\tfrac12)}{E_{x,t}},
  \label{eq:halfcount}
\end{equation}
so that both sides of every rate-scale score live on the same lattice
(Addendum~2~\S2, Addendum~3~\S5). Cells are never dropped for having no
deaths; the number of test cells in which the correction binds
($D_{x,t}<\tfrac12$, counted on the full age range before any model's age
mask) is reported per population, sex and regime. Such cells are
concentrated in the smallest populations, ISL and LUX, at the youngest and
oldest ages, where they are a material fraction of the 2020--2024 test
window; they are rare elsewhere. The same convention is applied to the
conformal calibration residuals (Section~\ref{sec:design-mechanisms}), where
a single zero-death calibration cell under a small rate floor would
otherwise inflate a band radius by an order of magnitude.

One deliberate exception: the \emph{observed} period life table used for
the realised $e_0$, $e_{65}$ and $\annuity$ of \Hfive\ keeps the raw
fractional deaths and a negligible rate floor, because at an age where
nobody died the crude rate $m_x=0$ is the correct life-table statement,
while a half death would shift the observed $e_0$ of a small population by
more than validation gate~3 tolerates. The two conventions---half a death
on the rate scale, no death in the observed life table, for the same
cell---are both registered (Addendum~3~\S5).

One historical detail of the triangle split matters for the placebo regime.
HMD's linear model for apportioning $1\times1$ death squares into Lexis
triangles (Protocol~v6, Appendix~A) carries a 1918--1919 influenza correction
extrapolated from Sweden and France. It changes how deaths are shared between
the two triangles of a square, not the $1\times1$ counts we use, and leaves
only a small model-based imprint on the 1918--1919 exposures.

\subsection{The likelihood}

Deaths are modelled as counts: $D_{x,t}\sim\mathrm{Poisson}(E_{x,t}\,m_{x,t})$
with $\log E_{x,t}$ as offset, following \citet{brouhns2002poisson}.
Nothing in the study models rates as Gaussian on the raw scale. The
classical families are fit by this likelihood, with two deliberate
exceptions retained as historical baselines: the original singular-value
Lee--Carter with its least-squares fit on log rates, and the CBD model with
its per-year least-squares fit on the logit of $q_{x,t}$. The neural
families fit the same Poisson deviance with the log-exposure offset
(neural-LC and the shallow CNN), a negative-binomial log-likelihood
(the distributional head), or, for the LSTM on $\kappa_t$, inherit the
SVD Lee--Carter stage and fit the index by squared error, as their
reference papers do; the multi-output Gaussian process places a Gaussian
likelihood on the half-count log rate~\eqref{eq:halfcount} with a
Poisson-level noise floor. The loss of every family is written in the
family specification committed before the neural sweep
(\texttt{docs/NEURAL-SPEC.md}) and summarised in
Section~\ref{sec:design-families}.

\subsection{Populations by regime}

Population codes are HMD's. The sets were fixed in the pre-registration from
an audit of the files on disk, not from results.

\paragraph{Shift (primary), 20 populations.} Those whose final annual data
reach 2024 in this vintage: BEL, CHE, CHL, DNK, EST, FIN, HKG, HRV, ISL, JPN,
KOR, LTU, LUX, LVA, NOR, PRT, SVK, SWE, TWN, USA. Every population trains on
the maximal contiguous block of complete years ending in 2019, from its own
first year (Addendum~3~\S2). The rule binds once: Belgium's 1914--1918
deaths and exposures are missing at every age---1{,}000 modelled cells,
$100$ ages $\times$ $2$ sexes $\times$ $5$ years, the only genuinely missing
data in the shift training panel (Addendum~3~\S2)---so BEL trains on
1919--2019; without the
rule a random-walk drift would treat 1913 to 1919 as one step and a
positional cohort index would mislabel every cohort after 1913 by five
years. The series lengths at the 2019 origin are very unequal, from
seventeen years (KOR) and nineteen (HRV) to more than two centuries (SWE);
this matters for which uncertainty mechanisms can be constructed on the
shortest panels (Section~\ref{sec:design-grid}).

\paragraph{Stable (control).} The same twenty populations, each from its own
first year of data, at expanding origins $T=1990,1992,\dots,2014$ with test
years $T+1,\dots,T+5$ capped at 2019. An (origin, population, sex) triple is
admissible only with at least fifteen complete training years
(Addendum~3~\S4): shorter windows produce plausible-looking junk, such as a
random-walk innovation variance estimated from two differences. The stable
panel is therefore unbalanced---the number of buildable population--sex
series grows across origins as the late-starting Baltic, East Asian and
South American series pass the floor---and the effective cluster count per
origin is reported alongside every stable-regime summary.

\paragraph{Placebo break, 11 populations.} Those with continuous coverage
over 1908--1922 and a start year no later than 1900: CHE, DNK, FIN, FRATNP,
GBRTENW, GBR\_SCO, ISL, ITA, NLD, NOR, SWE. Belgium is excluded because its
occupation-era data for 1914--1918 are missing at every age: 1{,}000
modelled cells in the registered unit of Addendum~3~\S2 ($100$ ages
$\times$ $2$ sexes $\times$ $5$ years), which appear in the raw $1\times1$
file as 555 rows (111 ages including $110{+}$, one row per age-year, 5
years) that are present but carry HMD's ``.'' missing marker in the Female,
Male and Total columns; HMD's own documentation warns of the gap. FRACNP and GBRCENW are excluded as
civilian-population variants overlapping FRATNP and GBRTENW; keeping both
would double-count one country, and HMD additionally flags FRACNP's 1919
denominators as implausible, so it is not used even as a sensitivity.

Addendum~1 to the pre-registration (2026-08-26, still before any real-data
fit) stratifies the eleven from HMD's country documentation into
\emph{neutral, register-based} series whose only break in the window is the
1918 influenza pandemic (CHE, DNK, FIN, ISL, NLD, NOR, SWE);
\emph{belligerent total-population} series in which HMD reconstructs
military deaths at male ages of roughly 18--45 (FRATNP, GBRTENW, ITA); and one
\emph{civilian-only} series in which deaths abroad are excluded by
construction (GBR\_SCO). The neutral stratum is the clean analogue for the
1918-versus-2020 comparison; the pooled placebo result is always reported
alongside the strata.

\begin{table}[htbp]\centering
\caption{Populations by regime: first and last year of single-age data in
the pinned vintage, training years available at the shift origin,
zero-exposure cells at ages $0$--$99$, and zero-death test cells
($D<\tfrac12$) per sex in the shift window. Generated from the data files,
not from any model output.}
\label{tab:populations}
\inputtable{tab-populations}
\end{table}

\subsection{Data-quality caveats carried into the analysis}

\begin{itemize}
\item \emph{Exposure conventions.} Exposures are estimated from census or
  register populations and births, with conventions that differ by country
  and era; \citet{cairns2016phantoms} show how cohort-size errors of several
  percent arise from uneven historical births. We take HMD exposures as
  given and note that any such error is common to every model in a cell.
\item \emph{The placebo window.} Pre-1950 data are patchier and the war years
  are, for some countries, reconstructed. The public HMD background and
  documentation file of each placebo population, the Methods Protocol and
  the release notes were read before any placebo result
  (\texttt{docs/DATA-PREREQS.md}); the per-country ``Notes'' files behind
  the HMD login were not consulted. No population carries an explicit
  1914--1922 quality flag, and the concerns that do exist are the strata
  above plus three specific items that become pre-specified sensitivities:
  GBR\_SCO's 1912--1938 population is interpolated by a spline HMD itself
  calls unreliable (sensitivity: drop GBR\_SCO); Denmark gained South Jutland
  in 1921, inside the test window (sensitivity: DNK with and without
  1921--1922); and the neutral stratum alone is reported as the cleanest
  pandemic-only comparison.
\item \emph{Revisability of 2024.} HMD carries no per-year provisional flag;
  each release re-issues whole series, so ``final through 2024'' means final
  as of this release. Known denominator risks in the shift panel: USA
  2020--2025 exposures blend preliminary 2020-census inputs; CHL exposures
  rest on the 2017 census pending revision; TWN 2024 deaths are provisional
  according to the short-term mortality fluctuations series. Register-based
  populations (the Nordic countries, CHE, EST, LVA, LTU) are effectively
  final. The pre-registered robustness run drops 2024 and re-scores the shift
  regime at $h\le4$; Addendum~1 adds dropping USA and CHL, a
  register-versus-census contrast, and a re-run on the next HMD vintage if
  one is released before submission.
\item \emph{Differing observation windows.} Series start between the
  mid-eighteenth century and the 2000s; the stable regime is therefore an
  unbalanced panel over origins, and the shortest shift panels cannot
  support every mechanism. This is why every inference is clustered by
  population rather than pooled over cells, and why every contrast is
  computed on the cells common to the arms it compares
  (Section~\ref{sec:design-grid}).
\end{itemize}

\subsection{Life-table conventions}

Life expectancy and annuity factors are computed from every predictive
sample's $m_x$ vector by one period-life-table routine:
$q_x = m_x/(1+(1-a_x)m_x)$; $a_0 = 0.07+1.7\,m_0$ capped to $[0.01,0.35]$,
a linear infant-separation rule in the spirit of \citet{andreev2015average};
$a_x=\tfrac12$ elsewhere; and constant-hazard closure at the last age,
$e_A = 1/m_A$. HMD itself uses the piecewise Andreev--Kingkade coefficients
and its own old-age treatment, so the two tables are not identical;
validation gate~3 requires our $e_0$ and $e_{65}$ on HMD's published $m_x$
to reproduce HMD's published values within 0.15 years, which is the size of
the documented $a_x$ simplifications. Annuity factors follow the standard
life-contingency definition \citep{dickson2020actuarial} on the period table
treated as static (Section~\ref{sec:metrics-actuarial}). A family defined
only on part of the age range (CBD, ages 55--99) yields a table truncated at
its first age: $e_0$ is then undefined and $e_{65}$ and $\annuity$ are
computed on the truncated table, on the predictive and observed sides
alike, with the age range recorded.
